% !TEX program = pdflatex
\documentclass[10pt,aspectratio=169]{beamer}

\usepackage[utf8]{inputenc}
\usepackage[T1]{fontenc}
\usepackage[english]{babel}
\usepackage{lmodern}
\usepackage{microtype}
\usepackage{amsmath,amssymb,bm,mathtools}
\usepackage{booktabs}
\usepackage{array}

\usetheme[numbering=fraction,progressbar=frametitle]{metropolis}
\setbeamertemplate{frame numbering}[fraction]
\setbeamertemplate{navigation symbols}{}

\definecolor{CourseBlue}{HTML}{1A3C68}
\definecolor{CourseGreen}{HTML}{216E39}
\definecolor{CourseGold}{HTML}{B8860B}
\definecolor{CourseRed}{HTML}{9A2C2C}
\definecolor{SoftBlue}{HTML}{F2F6FA}
\definecolor{SoftGreen}{HTML}{F1F8F3}
\definecolor{SoftRed}{HTML}{FFF3F3}

\setbeamercolor{normal text}{fg=black,bg=white}
\setbeamercolor{structure}{fg=CourseBlue}
\setbeamercolor{title}{fg=CourseBlue}
\setbeamerfont{title}{size=\Large,series=\bfseries}
\setbeamercolor{frametitle}{bg=CourseBlue,fg=white}
\setbeamerfont{frametitle}{series=\bfseries,size=\large}
\setbeamertemplate{frametitle}{
  \nointerlineskip
  \begin{beamercolorbox}[wd=\paperwidth,ht=3.2ex,dp=1.4ex,leftskip=1em,rightskip=1em]{frametitle}
    \usebeamerfont{frametitle}\insertframetitle
  \end{beamercolorbox}
}
\setbeamercolor{block title}{bg=CourseBlue,fg=white}
\setbeamercolor{block body}{bg=SoftBlue,fg=black}
\setbeamercolor{block title example}{bg=CourseGreen,fg=white}
\setbeamercolor{block body example}{bg=SoftGreen,fg=black}
\setbeamercolor{block title alerted}{bg=CourseRed,fg=white}
\setbeamercolor{block body alerted}{bg=SoftRed,fg=black}
\setbeamertemplate{blocks}[rounded][shadow=false]

\newenvironment{modelbox}{%
  \begin{beamercolorbox}[rounded=true,sep=0.7em,wd=\textwidth]{block body}
  \textbf{\color{CourseBlue}MODEL ASSUMPTION}\quad
}{\end{beamercolorbox}}
\newenvironment{marketbox}{%
  \begin{beamercolorbox}[rounded=true,sep=0.7em,wd=\textwidth]{block body example}
  \textbf{\color{CourseGreen}MARKET ASSUMPTION}\quad
}{\end{beamercolorbox}}
\newenvironment{simplificationbox}{%
  \begin{beamercolorbox}[rounded=true,sep=0.7em,wd=\textwidth]{block body alerted}
  \textbf{\color{CourseRed}SIMPLIFICATION}\quad
}{\end{beamercolorbox}}
\newenvironment{pausebox}{%
  \begin{beamercolorbox}[rounded=true,sep=0.8em,wd=\textwidth]{block body}
  \textbf{\color{CourseGold}YOUR TURN}\quad
}{\end{beamercolorbox}}

\newcommand{\dd}{\,\mathrm d}
\newcommand{\E}{\mathbb E}
\newcommand{\Pp}{\mathbb P}
\newcommand{\Q}{\mathbb Q}
\newcommand{\R}{\mathbb R}
\newcommand{\cF}{\mathcal F}
\newcommand{\Ncdf}{\mathrm N}
\newcommand{\phiN}{\varphi}
\newcommand{\1}{\mathbf 1}
\DeclareMathOperator{\Var}{Var}
\DeclareMathOperator{\Cov}{Cov}

\AtBeginSection[]{
  \begin{frame}[plain]
    \vfill
    {\Large\bfseries\color{CourseBlue}\insertsectionhead\par}
    \medskip
    {\color{CourseBlue}\rule{0.92\textwidth}{0.8pt}}
    \vfill
  \end{frame}
}

\title{Black--Scholes from Replication}
\subtitle{It\^o practice, self-financing strategies, the pricing PDE, verification, and Greeks}
\author{}
\date{Lecture 5}

\begin{document}

\begin{frame}[plain]
  \vfill
  {\LARGE\bfseries\color{CourseBlue}Black--Scholes from Replication\par}
  \medskip
  {\large It\^o practice, self-financing strategies, the pricing PDE,
  verification, and Greeks\par}
  \medskip
  {\color{CourseGold}\rule{0.94\textwidth}{0.8pt}}
  \medskip
  {Lecture 5\par}
  \vfill
\end{frame}

\begin{frame}{Today one idea will connect calculus, pricing, and hedging}
\begin{enumerate}
  \item Use It\^o's formula on three new objects.
  \item Define a self-financing trading strategy precisely.
  \item Replicate \(g(S_T)\) and derive the Black--Scholes PDE.
  \item Solve the PDE through time reversal, discounting, and \(\log S\).
  \item Verify the candidate price and recover the hedge.
  \item Specialize to a call and compute its Greeks.
\end{enumerate}
\begin{block}{Central message}
The price is the initial wealth of a self-financing strategy that reproduces
the payoff in every state of the model.
\end{block}
\end{frame}

% ============================================================
\section{Three New It\^o Calculations}
% ============================================================

\begin{frame}{The formula is a chain rule plus quadratic variation}
If
\[
\dd X_t=b_t\dd t+a_t\dd W_t,
\qquad f\in C^{1,2},
\]
then
\[
\boxed{
\dd f(t,X_t)
=\left(f_t+b_tf_x+\frac12a_t^2f_{xx}\right)\dd t
+a_tf_x\dd W_t.}
\]
\begin{columns}[T,onlytextwidth]
\column{0.48\textwidth}
\begin{block}{Calculation habit}
Compute \(f_t,f_x,f_{xx}\) first. Substitute the process only afterward.
\end{block}
\column{0.48\textwidth}
\begin{alertblock}{Logical habit}
Zero drift first gives a local martingale. Check integrability before claiming
a true martingale.
\end{alertblock}
\end{columns}
\end{frame}

\begin{frame}{A two-dimensional Brownian motion gives drift 2}
Let \(W=(W^1,W^2)\) be a standard two-dimensional Brownian motion and
\[
R_t=\|W_t\|^2=(W_t^1)^2+(W_t^2)^2.
\]
This is the squared radius, or squared Bessel process of dimension 2.
\begin{pausebox}
Compute \(\dd R_t\). Which deterministic term makes the result a martingale?
\end{pausebox}
\pause
Applying It\^o to each squared coordinate gives
\[
\boxed{\dd R_t=2W_t^\top\dd W_t+2\dd t,}
\qquad
\boxed{R_t-2t=2\int_0^tW_s^\top\dd W_s.}
\]
On bounded horizons the integrand is square-integrable, so \(R_t-2t\) is a
true martingale.
\end{frame}

\begin{frame}{In d dimensions, the drift counts independent directions}
Now let \(W\) be a standard \(d\)-dimensional Brownian motion. Then
\[
\E\|W_t\|^2=dt,
\qquad
[\,\|W\|^2\,]_t=4\int_0^t\|W_s\|^2\dd s.
\]
\begin{exampleblock}{Interpretation}
Each Brownian coordinate contributes one unit of quadratic variation per unit
of time. Dimension appears as drift after a nonlinear transformation.
\end{exampleblock}
\begin{simplificationbox}
The coordinates are independent and normalized:
\(\dd[W^i,W^j]_t=\delta_{ij}\dd t\). With covariance matrix \(\Sigma\), the
correction becomes \(\operatorname{tr}(\Sigma)\dd t\).
\end{simplificationbox}
\end{frame}

\begin{frame}{Example 2 combines a finite-variation state with Brownian motion}
Define integrated Brownian motion
\[
A_t=\int_0^tW_s\dd s,
\qquad Y_t=W_tA_t.
\]
\begin{pausebox}
Use the product rule to find \(\dd Y_t\), then compute \(\E[Y_t]\).
\end{pausebox}
\pause
Since \(\dd A_t=W_t\dd t\) and \(A\) has finite variation,
\(\dd[W,A]_t=0\). Therefore
\[
\boxed{\dd(W_tA_t)=A_t\dd W_t+W_t^2\dd t.}
\]
Taking expectations,
\[
\boxed{\E[W_tA_t]=\int_0^t\E[W_s^2]\dd s=\frac{t^2}{2}.}
\]
\end{frame}

\begin{frame}{The same answer is a covariance computed two ways}
Because both variables are centered,
\[
\E[W_tA_t]=\Cov\left(W_t,\int_0^tW_s\dd s\right).
\]
\begin{pausebox}
Verify the answer directly from \(\Cov(W_t,W_s)=\min(t,s)\).
\end{pausebox}
\pause
\[
\Cov\left(W_t,\int_0^tW_s\dd s\right)
=\int_0^t\Cov(W_t,W_s)\dd s
=\int_0^t s\dd s
=\frac{t^2}{2}.
\]
\begin{block}{Why this example matters}
It\^o calculus handles path-dependent state variables after we enlarge the
state from \(W_t\) to \((W_t,A_t)\).
\end{block}
\end{frame}

\begin{frame}{Example 3 turns a future event into a martingale}
Fix \(T>0\). For \(t<T\), consider
\[
u(t,x)=\Ncdf\!\left(\frac{x}{\sqrt{T-t}}\right),
\]
where \(\Ncdf\) and \(\phiN\) are the standard normal cdf and density.
\begin{pausebox}
Compute \(u_t,u_x,u_{xx}\). Does \(u(t,W_t)\) have drift?
\end{pausebox}
\pause
With \(\tau=T-t\) and \(z=x/\sqrt{\tau}\),
\[
u_x=\frac{\phiN(z)}{\sqrt{\tau}},
\qquad
u_{xx}=-\frac{x\phiN(z)}{\tau^{3/2}},
\qquad
u_t=\frac{x\phiN(z)}{2\tau^{3/2}}.
\]
\end{frame}

\begin{frame}{The backward heat equation removes the drift}
The derivatives satisfy
\[
\boxed{u_t+\frac12u_{xx}=0.}
\]
It\^o's formula therefore gives
\[
\boxed{
\dd u(t,W_t)
=\frac{\phiN(W_t/\sqrt{T-t})}{\sqrt{T-t}}\,\dd W_t.}
\]
\pause
Moreover,
\[
u(t,W_t)=\Pp(W_T>0\mid\cF_t),
\]
because \(W_T=W_t+\sqrt{T-t}\,Z\), with \(Z\sim N(0,1)\) independent of
\(\cF_t\).
\begin{exampleblock}{Preview}
A terminal payoff, a backward PDE, and a conditional expectation are three
views of the same object. Black--Scholes will add discounting and replication.
\end{exampleblock}
\end{frame}


% ============================================================
\section{Replication Creates the Pricing Equation}
% ============================================================

\begin{frame}{Black--Scholes starts with a model and a trading technology}
\begin{columns}[T,onlytextwidth]
\column{0.49\textwidth}
\begin{modelbox}
One stock follows
\[
\frac{\dd S_t}{S_t}=\mu\dd t+\sigma\dd W_t,\quad \sigma>0.
\]
One bank account satisfies
\[
\dd B_t=rB_t\dd t,\quad B_t=e^{rt}.
\]
\end{modelbox}
\column{0.49\textwidth}
\begin{marketbox}
Trading is continuous and frictionless. Stock and bank account can be bought
or shorted, and borrowing/lending occurs at the same rate \(r\).
\end{marketbox}
\end{columns}
\medskip
\begin{simplificationbox}
For now, \(r,\mu,\sigma\) are constant, the stock pays no dividend, and its
paths are continuous.
\end{simplificationbox}
\end{frame}

\begin{frame}{The assumptions have different logical roles}
\small
\begin{tabular}{>{\bfseries}p{0.18\textwidth}p{0.34\textwidth}p{0.38\textwidth}}
\toprule
Type & Assumption & Where it enters\\
\midrule
Model & GBM, constant \(\sigma\), no jumps & It\^o dynamics and PDE coefficients\\
Market & No arbitrage & A locally riskless position must earn \(r\)\\
Market & Continuous frictionless trading & Delta can be adjusted continuously\\
Market & One Brownian risk, one stock & Every smooth European claim can be hedged\\
Simplification & Constant \(r,\sigma\), no dividends & Closed-form \(d_\pm\) formula\\
Regularity & \(p\in C^{1,2}\) before maturity & Classical It\^o formula applies\\
\bottomrule
\end{tabular}
\begin{alertblock}{Do not merge these claims}
A realistic price model need not create a complete market; a no-arbitrage
market need not admit a closed-form formula.
\end{alertblock}
\end{frame}

\begin{frame}{A European claim fixes only the terminal value}
The claim pays
\[
H=g(S_T)
\]
at maturity \(T\). We seek a deterministic Markov price
\[
V_t=p(t,S_t),
\qquad
\boxed{p(T,s)=g(s).}
\]
\begin{modelbox}
The Markov representation assumes the current stock price is a sufficient
state variable for this claim.
\end{modelbox}
\begin{simplificationbox}
We initially assume \(p\in C^{1,2}([0,T)\times(0,\infty))\). Vanilla payoffs
have a kink at maturity, but the price becomes smooth for \(t<T\).
\end{simplificationbox}
\end{frame}

\begin{frame}{A trading strategy separates holdings from gains}
At time \(t\), hold
\[
\Delta_t\ \text{shares of stock},
\qquad
\beta_t\ \text{units of the bank account}.
\]
The portfolio value is
\[
\boxed{X_t=\Delta_tS_t+\beta_tB_t.}
\]
\begin{marketbox}
\(\Delta\) and \(\beta\) must be predictable: the position used over the next
instant is chosen using currently available information.
\end{marketbox}
\begin{block}{Replication objective}
Choose \(X_0,\Delta,\beta\) so that \(X_T=g(S_T)\) in every state allowed by
the model.
\end{block}
\end{frame}

\begin{frame}{Self-financing is a budget constraint on rebalancing}
\begin{center}
\Large
\[
\text{change in wealth}
=
\text{gain on the assets}
+\text{external cash flow}.
\]
\end{center}
\pause
\begin{block}{Self-financing}
After the initial investment, the external cash flow is zero.
\end{block}
\begin{columns}[T,onlytextwidth]
\column{0.48\textwidth}
\begin{exampleblock}{Allowed}
\(\Delta\) and \(\beta\) may change through time as information arrives.
\end{exampleblock}
\column{0.48\textwidth}
\begin{alertblock}{Not allowed}
Rebalancing cannot create wealth or require a new deposit.
\end{alertblock}
\end{columns}
\medskip
We first write this budget constraint in discrete time, where every cash flow
is visible, and then pass to continuous time.
\end{frame}

\begin{frame}{In discrete time, rebalancing preserves current wealth}
\small
At date \(t_n\), value the old and new holdings at the \emph{same} prices:
\[
X_n^-=\Delta_{n-1}S_n+\beta_{n-1}B_n,
\qquad
X_n=\Delta_nS_n+\beta_nB_n.
\]
\begin{pausebox}
If no cash enters or leaves at \(t_n\), what equality must hold?
\end{pausebox}
\pause
\[
X_n=X_n^-
\quad\Longleftrightarrow\quad
\boxed{
S_n(\Delta_n-\Delta_{n-1})
+B_n(\beta_n-\beta_{n-1})=0.}
\]
\begin{block}{Interpretation}
Buying additional shares is financed by reducing cash; selling shares
increases cash. Rebalancing only changes the composition of wealth.
\end{block}
\end{frame}

\begin{frame}{In discrete time, gains come only from asset-price moves}
\small
After rebalancing at \(t_n\), keep \((\Delta_n,\beta_n)\) over
\((t_n,t_{n+1}]\). Just before the next rebalance,
\[
X_{n+1}^-=\Delta_nS_{n+1}+\beta_nB_{n+1}.
\]
\begin{pausebox}
Subtract \(X_n=\Delta_nS_n+\beta_nB_n\). What is the wealth increment?
\end{pausebox}
\pause
\[
\boxed{
X_{n+1}^--X_n
=\Delta_n(S_{n+1}-S_n)
+\beta_n(B_{n+1}-B_n).}
\]
\begin{exampleblock}{Key distinction}
Rebalancing changes holdings but not wealth. Asset-price changes generate
gains and losses.
\end{exampleblock}
\end{frame}

\begin{frame}{Continuous self-financing is the limit of discrete gains}
\small
On a trading grid, accumulated wealth satisfies
\[
X_t-X_0
=\sum_n\Delta_n(S_{n+1}-S_n)
+\sum_n\beta_n(B_{n+1}-B_n).
\]
As the mesh tends to zero, the predictable left-point sums become stochastic
integrals:
\[
\boxed{
X_t=X_0+\int_0^t\Delta_u\dd S_u+\int_0^t\beta_u\dd B_u,}
\qquad
\boxed{\dd X_t=\Delta_t\dd S_t+\beta_t\dd B_t.}
\]
\pause
Together with \(X_t=\Delta_tS_t+\beta_tB_t\), the product rule yields the
equivalent rebalancing constraint
\[
\boxed{
S_t\dd\Delta_t+\dd[\Delta,S]_t+B_t\dd\beta_t=0.}
\]
\end{frame}

\begin{frame}{First calculation: expand the option dynamics}
\[
V_t=p(t,S_t),
\qquad
\dd S_t=\mu S_t\dd t+\sigma S_t\dd W_t.
\]
\begin{pausebox}
Use It\^o's formula. Separate the drift and diffusion coefficients.
\end{pausebox}
\pause
\[
\boxed{
\begin{aligned}
\dd V_t
={}&\left(
p_t+\mu S_tp_s+\frac12\sigma^2S_t^2p_{ss}
\right)\dd t\\
&+\sigma S_tp_s\dd W_t.
\end{aligned}}
\]
All derivatives are evaluated at \((t,S_t)\).
\end{frame}

\begin{frame}{The bank-account gain comes from the cash balance}
\small
The portfolio value is
\[
X_t=\Delta_tS_t+\beta_tB_t
\quad\Longrightarrow\quad
\boxed{\beta_tB_t=X_t-\Delta_tS_t.}
\]
Since \(\dd B_t=rB_t\dd t\), the bank-account gain is
\[
\boxed{\beta_t\dd B_t
=r\beta_tB_t\dd t
=r(X_t-\Delta_tS_t)\dd t.}
\]
\begin{pausebox}
Now insert \(\dd S_t=\mu S_t\dd t+\sigma S_t\dd W_t\).
\end{pausebox}
\pause
\[
\boxed{
\dd X_t
=\left[rX_t+(\mu-r)S_t\Delta_t\right]\dd t
+\sigma S_t\Delta_t\dd W_t.}
\]
\begin{marketbox}
Here \(r\) applies to both positive and negative bank balances.
\end{marketbox}
\end{frame}

\begin{frame}{Matching random shocks determines the hedge}
Replication requires
\[
X_t=V_t=p(t,S_t)
\quad\text{and}\quad
\dd X_t=\dd V_t.
\]
\begin{pausebox}
Match the \(\dd W_t\) coefficients. What must \(\Delta_t\) be?
\end{pausebox}
\pause
\[
\sigma S_t\Delta_t=\sigma S_tp_s(t,S_t).
\]
Since \(S_t>0\) and \(\sigma>0\),
\[
\boxed{\Delta_t=p_s(t,S_t).}
\]
\begin{block}{Economic meaning}
Delta is not merely a sensitivity. In this complete diffusion model, it is
the stock holding in the replicating strategy.
\end{block}
\end{frame}

\begin{frame}{Matching drifts produces one final cancellation}
Set \(X_t=p(t,S_t)\) and \(\Delta_t=p_s(t,S_t)\).
\begin{pausebox}
Match the drift coefficients and simplify. Which parameter disappears?
\end{pausebox}
\pause
\[
rp+(\mu-r)sp_s
=p_t+\mu sp_s+\frac12\sigma^2s^2p_{ss}.
\]
The terms containing the physical drift cancel:
\[
\boxed{
p_t+rsp_s+\frac12\sigma^2s^2p_{ss}-rp=0,}
\qquad
\boxed{p(T,s)=g(s).}
\]
\end{frame}

\begin{frame}{The Black--Scholes PDE is a replication equation}
\[
\boxed{
\begin{cases}
p_t(t,s)+rs\,p_s(t,s)
+\dfrac12\sigma^2s^2p_{ss}(t,s)-rp(t,s)=0,\\[0.4em]
p(T,s)=g(s).
\end{cases}}
\]
\begin{exampleblock}{What created each term?}
\begin{itemize}
  \item \(p_t\): passage of calendar time;
  \item \(rsp_s\): financing the stock position;
  \item \(\tfrac12\sigma^2s^2p_{ss}\): quadratic variation and gamma;
  \item \(-rp\): financing the total option value.
\end{itemize}
\end{exampleblock}
\end{frame}

\begin{frame}{Why does the expected stock return disappear?}
The stock exposure appears twice:
\[
\underbrace{\mu sp_s}_{\text{option's stock exposure}}
\qquad\text{and}\qquad
\underbrace{\mu s\Delta}_{\text{replicating stock position}}.
\]
With \(\Delta=p_s\), they are identical and cancel.
\begin{block}{Pricing conclusion}
In the model, valuation depends on \(r\) and \(\sigma\), not on the estimated
physical drift \(\mu\).
\end{block}
\begin{alertblock}{What this does not say}
Real-world probabilities remain essential for forecasting, risk management,
and expected returns.
\end{alertblock}
\end{frame}

\begin{frame}{Replication needs no arbitrage and enough traded risk}
\begin{columns}[T,onlytextwidth]
\column{0.48\textwidth}
\begin{marketbox}
\textbf{No arbitrage.} Two self-financing portfolios with the same terminal
payoff must have the same initial value.
\end{marketbox}
\column{0.48\textwidth}
\begin{marketbox}
\textbf{Completeness.} One stock spans the one Brownian source of uncertainty,
so the option's diffusion exposure can be matched.
\end{marketbox}
\end{columns}
\medskip
\begin{alertblock}{When the argument breaks}
With jumps, stochastic volatility driven by another shock, trading
constraints, or transaction costs, one stock generally cannot eliminate all
risk. The replication price may cease to be unique.
\end{alertblock}
\end{frame}

\begin{frame}{The bank-account holding follows from the value identity}
Once \(\Delta_t=p_s(t,S_t)\),
\[
\beta_t
=\frac{p(t,S_t)-S_tp_s(t,S_t)}{B_t}.
\]
\begin{pausebox}
At maturity, what does the portfolio hold when \(g(s)=(s-K)^+\), away from
the strike?
\end{pausebox}
\pause
\[
\Delta_T=
\begin{cases}1,&S_T>K,\\0,&S_T<K,\end{cases}
\qquad
\beta_TB_T=
\begin{cases}-K,&S_T>K,\\0,&S_T<K.\end{cases}
\]
The call is replicated by one share financed by \(K\) of borrowing when it
finishes in the money.
\end{frame}

\begin{frame}{Self-financing is not the same as zero transaction costs}
\small
\begin{marketbox}
The model has no commissions or bid--ask spread. But buying shares still
requires cash: the bank-account holding must adjust to pay for them.
\end{marketbox}
\begin{enumerate}
\setlength{\itemsep}{0.2em}
  \item \textbf{Rebalancing does not create wealth.} Buying shares is paid
        for from the cash account; this is funding, not a transaction fee.
  \item \textbf{Matching Delta removes only the Brownian shock.} The
        portfolio must also satisfy the self-financing wealth equation.
  \item \textbf{The option need not earn \(r\) by itself.} After matching
        its stock exposure, the whole locally riskless portfolio earns \(r\).
  \item \textbf{\(\mu\) is not replaced by \(r\).} It cancels because the
        option and its replicating stock position have the same exposure.
\end{enumerate}
\end{frame}

% ============================================================
\section{Solve the PDE by Taking Logs}
% ============================================================

\begin{frame}{The terminal-value problem is our starting point}
\[
p_t+rsp_s+\frac12\sigma^2s^2p_{ss}-rp=0,
\qquad
p(T,s)=g(s).
\]
The coefficient \(s^2\) prevents this from being the standard heat equation.
\begin{block}{Four transformations}
\[
\text{time to maturity}
\longrightarrow
\text{remove discounting}
\longrightarrow
\log s
\longrightarrow
\text{remove transport}.
\]
\end{block}
\begin{simplificationbox}
Constant \(r\) and \(\sigma\) make these transformations explicit. The pricing
principle survives more general coefficients, but the closed form usually does
not.
\end{simplificationbox}
\end{frame}

\begin{frame}{Step 1 turns the terminal condition into initial data}
Set
\[
\tau=T-t,
\qquad
v(\tau,s)=p(T-\tau,s).
\]
\begin{pausebox}
Transform the PDE and its terminal condition.
\end{pausebox}
\pause
Because \(p_t=-v_\tau\),
\[
\boxed{
v_\tau
=rsv_s+\frac12\sigma^2s^2v_{ss}-rv,
\qquad
v(0,s)=g(s).}
\]
\end{frame}

\begin{frame}{Step 2 removes the financing term}
Define
\[
u(\tau,s)=e^{r\tau}v(\tau,s).
\]
\begin{pausebox}
Compute \(u_\tau,u_s,u_{ss}\) and simplify the PDE.
\end{pausebox}
\pause
\[
u_\tau=ru+e^{r\tau}v_\tau,
\qquad
u_s=e^{r\tau}v_s,
\qquad
u_{ss}=e^{r\tau}v_{ss}.
\]
Thus
\[
\boxed{
u_\tau=rsu_s+\frac12\sigma^2s^2u_{ss},
\qquad u(0,s)=g(s).}
\]
\end{frame}

\begin{frame}{Step 3 makes multiplicative motion additive}
Let
\[
x=\log s,
\qquad
w(\tau,x)=u(\tau,e^x).
\]
\begin{pausebox}
Express \(u_s\) and \(u_{ss}\) using \(w_x,w_{xx}\).
\end{pausebox}
\pause
\[
w_x=su_s,
\qquad
w_{xx}=su_s+s^2u_{ss}.
\]
Therefore
\[
\boxed{
u_s=\frac{w_x}{s},
\qquad
u_{ss}=\frac{w_{xx}-w_x}{s^2}.}
\]
\end{frame}

\begin{frame}{The log transform leaves heat plus constant transport}
Insert the derivatives:
\[
\boxed{
w_\tau
=\frac12\sigma^2w_{xx}
+\left(r-\frac12\sigma^2\right)w_x,
\qquad
w(0,x)=g(e^x).}
\]
\begin{pausebox}
Why does \(r-\sigma^2/2\) look familiar?
\end{pausebox}
\pause
It is the drift of \(\log S\) when the stock's pricing drift is \(r\):
\[
\dd\log S_t
=\left(r-\frac12\sigma^2\right)\dd t+\sigma\dd W_t.
\]
\end{frame}

\begin{frame}{Step 4 moves with the transport velocity}
Write
\[
a=r-\frac12\sigma^2,
\qquad
\varphi(\tau,y)=w(\tau,y-a\tau).
\]
\begin{pausebox}
Differentiate \(\varphi\). Which PDE does it satisfy?
\end{pausebox}
\pause
\[
\varphi_\tau=w_\tau-aw_x,
\qquad
\varphi_y=w_x,
\qquad
\varphi_{yy}=w_{xx}.
\]
Hence
\[
\boxed{
\varphi_\tau=\frac12\sigma^2\varphi_{yy},
\qquad
\varphi(0,y)=g(e^y).}
\]
\end{frame}

\begin{frame}{The heat equation averages terminal values with a Gaussian}
If \(Z\sim N(0,1)\), then
\[
\boxed{
\varphi(\tau,y)
=\E\!\left[g\!\left(e^{y+\sigma\sqrt{\tau}Z}\right)\right].}
\]
Equivalently,
\[
\varphi(\tau,y)
=\int_{\R}g(e^z)
\frac{\exp\!\left(-\frac{(z-y)^2}{2\sigma^2\tau}\right)}
{\sqrt{2\pi\sigma^2\tau}}\dd z.
\]
\begin{block}{Meaning}
Diffusion smooths the terminal payoff by averaging it over future Gaussian
log-price shocks.
\end{block}
\end{frame}

\begin{frame}{Undoing the transformations gives the general price}
For \(\tau=T-t\),
\[
\boxed{
\begin{aligned}
p(t,s)
=e^{-r\tau}\E\Big[
g\Big(s\exp\{&(r-\tfrac12\sigma^2)\tau\\
&+\sigma\sqrt{\tau}Z\}\Big)
\Big],
\qquad Z\sim N(0,1).
\end{aligned}}
\]
\begin{modelbox}
Inside the expectation, the auxiliary stock grows at drift \(r\), not at the
physical drift \(\mu\).
\end{modelbox}
\begin{alertblock}{Scope}
This is the general European-payoff formula in the Black--Scholes model. It is
a Gaussian integral, not a closed form for every \(g\).
\end{alertblock}
\end{frame}

\begin{frame}{Regularity and growth conditions are doing real work}
The formal derivation is valid for a smooth payoff of controlled growth.
For common nonsmooth payoffs:
\begin{itemize}
  \item the terminal condition may have kinks or discontinuities;
  \item the Gaussian convolution is smooth for every \(t<T\);
  \item the terminal condition is recovered in the appropriate limiting sense.
\end{itemize}
\begin{simplificationbox}
We suppress the distinction between classical, weak, and viscosity solutions.
For vanilla calls and puts, the pricing formula is well defined and the
required derivatives exist before maturity.
\end{simplificationbox}
\end{frame}

% ============================================================
\section{Feynman--Kac Verification}
% ============================================================

\begin{frame}{Feynman--Kac represents the PDE solution as an expectation}
\small
We want to turn the PDE solution \(p\) into an expectation. Its spatial
drift term is \(rs\,p_s\), and its discount term is \(-rp\).
\begin{pausebox}
Choose a process whose drift makes It\^o's formula reproduce those terms.
\end{pausebox}
\pause
For each starting pair \((t,s)\), define \(S^r\) on \([t,T]\) by
\[
\frac{\dd S_u^r}{S_u^r}=r\dd u+\sigma\dd W_u,
\qquad S_t^r=s.
\]
Write \(\Pp^r_{t,s}\) and \(\E^r_{t,s}\) for probability and expectation
under this process started at \(s\) at time \(t\).
\begin{simplificationbox}
\(S^r\) is an auxiliary pricing process, not the physical stock forecast.
Its drift \(r\) matches the PDE; later, a change of measure gives it a
risk-neutral interpretation.
\end{simplificationbox}
\end{frame}

\begin{frame}{Discounting and the PDE remove the drift}
\small
\begin{pausebox}
Apply It\^o's formula to \(e^{-r(u-t)}p(u,S_u^r)\). Identify its drift
and explain why that drift vanishes.
\end{pausebox}
\pause
It\^o's formula gives
\[
\begin{aligned}
\dd\!\left(e^{-r(u-t)}p(u,S_u^r)\right)
=e^{-r(u-t)}
\bigg[&p_t+rS_u^{r}p_s\\
&+\frac12\sigma^2(S_u^{r})^2p_{ss}-rp\bigg]\dd u\\
&+e^{-r(u-t)}\sigma S_u^{r}p_s\dd W_u.
\end{aligned}
\]
The drift bracket is exactly the Black--Scholes PDE, evaluated at
\((u,S_u^r)\), so it vanishes:
\[
\boxed{
\dd\!\left(e^{-r(u-t)}p(u,S_u^r)\right)
=e^{-r(u-t)}\sigma S_u^rp_s\dd W_u.}
\]
\end{frame}

\begin{frame}{The terminal condition turns the martingale into the price}
If the stochastic integral is a true martingale, its expected future value
equals its current value. Thus
\[
p(t,s)
=\E^r_{t,s}\!\left[
e^{-r(T-t)}p(T,S_T^{r})
\right].
\]
\pause
Now use the terminal condition \(p(T,x)=g(x)\):
\[
\boxed{
p(t,s)
=\E^r_{t,s}\!\left[
e^{-r(T-t)}g(S_T^{r})
\right].}
\]
\begin{alertblock}{Verification condition}
For example, it is enough that the displayed stochastic integral has finite
expected quadratic variation. Otherwise, localization and uniform
integrability are needed.
\end{alertblock}
\end{frame}

\begin{frame}{The same price function gives the hedge in the actual market}
\small
Now return to the traded stock \(S_t\), whose physical drift is \(\mu\).
Given a sufficiently regular PDE solution, define
\[
\Delta_t=p_s(t,S_t),
\qquad
\beta_t=\frac{p(t,S_t)-S_tp_s(t,S_t)}{B_t}.
\]
\begin{pausebox}
Check that these holdings are self-financing and match the payoff.
\end{pausebox}
\pause
The PDE rearranges It\^o's formula as
\[
\dd p(t,S_t)
=p_s(t,S_t)\dd S_t
+r\bigl(p(t,S_t)-S_tp_s(t,S_t)\bigr)\dd t.
\]
Therefore
\[
\dd p(t,S_t)=\Delta_t\dd S_t+\beta_t\dd B_t,
\qquad
p(T,S_T)=g(S_T).
\]
\begin{block}{Two roles, one price function}
\(S^r\) represents the price as an expectation; \(S\) is the traded asset
used to implement the replicating hedge.
\end{block}
\end{frame}

\begin{frame}{Expectation and replication prove different facts}
\begin{columns}[T,onlytextwidth]
\column{0.48\textwidth}
\begin{block}{PDE / replication}
Explains why the price is unique, why \(\mu\) disappears, and why
\(\Delta=p_s\).
\end{block}
\column{0.48\textwidth}
\begin{exampleblock}{Expectation / verification}
Computes the PDE solution and checks terminal matching and discounted
martingale dynamics.
\end{exampleblock}
\end{columns}
\medskip
\[
\boxed{
\text{self-financing replication}
\Longleftrightarrow
\text{PDE}
\Longleftrightarrow
\text{discounted expectation}.}
\]
\begin{simplificationbox}
We built an auxiliary drift-\(r\) process directly. A later measure-change
argument identifies it with risk-neutral pricing under \(\Q\).
\end{simplificationbox}
\end{frame}

% ============================================================
\section{The European Call Formula}
% ============================================================

\begin{frame}{The call payoff splits into two truncated expectations}
Let
\[
g(s)=(s-K)^+,
\qquad \tau=T-t.
\]
Under the auxiliary drift-\(r\) dynamics,
\[
S_T^r
=s\exp\left(
(r-\tfrac12\sigma^2)\tau+\sigma\sqrt{\tau}Z
\right).
\]
Therefore
\[
C(t,s)
=e^{-r\tau}\left(
\E^r_{t,s}[S_T^r\1_{\{S_T^r>K\}}]
-K\Pp^r_{t,s}(S_T^r>K)
\right).
\]
\end{frame}

\begin{frame}{The exercise event creates \(d_-\)}
\begin{pausebox}
Rewrite \(S_T^r>K\) as a threshold event for \(Z\sim N(0,1)\).
\end{pausebox}
\pause
\[
S_T^r>K
\iff
Z>
\frac{\log(K/s)-(r-\sigma^2/2)\tau}{\sigma\sqrt{\tau}}
=-d_-,
\]
where
\[
\boxed{
d_-=
\frac{\log(s/K)+(r-\tfrac12\sigma^2)\tau}
{\sigma\sqrt{\tau}}.}
\]
Thus
\[
\boxed{\Pp^r_{t,s}(S_T^r>K)=\Ncdf(d_-).}
\]
\end{frame}

\begin{frame}{A stock-weighted exercise probability shifts the Gaussian}
\small
The remaining call term is
\(\E^r_{t,s}[e^{aZ}\1_{\{Z>-d_-\}}]\), where \(a=\sigma\sqrt{\tau}\).
Unlike the strike term, it weights exercise states by the stock price and
therefore favors larger \(Z\).
\begin{pausebox}
Complete the square to evaluate this weighted Gaussian tail.
\end{pausebox}
\pause
Letting $\varphi$ be the standard gaussian density. 
\[
\begin{aligned}
e^{az}\phiN(z)&=e^{a^2/2}\phiN(z-a),\\
\E^r_{t,s}[e^{aZ}\1_{\{Z>-d_-\}}]
&=e^{a^2/2}\Ncdf(d_-+a).
\end{aligned}
\]
Thus \(\boxed{d_+=d_-+a}\): \(d_-\) measures exercise probability, while
\(d_+\) prices the stock delivered in those exercise states.
\end{frame}

\begin{frame}{The Black--Scholes call formula follows}
The stock part becomes
\[
\E^r_{t,s}[S_T^r\1_{\{S_T^r>K\}}]
=se^{r\tau}\Ncdf(d_+).
\]
Hence
\[
\boxed{
C(t,s)=s\Ncdf(d_+)-Ke^{-r\tau}\Ncdf(d_-),}
\]
with
\[
\boxed{
d_\pm
=\frac{\log(s/K)+r\tau}{\sigma\sqrt{\tau}}
\pm\frac12\sigma\sqrt{\tau}.}
\]
\begin{simplificationbox}
This closed form uses constant \(r,\sigma\), no dividends, European exercise,
and lognormal stock dynamics.
\end{simplificationbox}
\end{frame}

\begin{frame}{The two normal terms have different meanings}
\[
C(t,s)
=\underbrace{s\Ncdf(d_+)}_{\text{stock delivered in exercise states}}
-\underbrace{Ke^{-r\tau}\Ncdf(d_-)}_{\text{discounted strike paid}}.
\]
\begin{alertblock}{Probability warning}
\(\Ncdf(d_-)\) is the exercise probability under the auxiliary pricing
dynamics. \(\Ncdf(d_+)\) is not the same probability; it arises from a
stock-weighted expectation.
\end{alertblock}
\begin{block}{Useful relation}
\[
d_+-d_-=\sigma\sqrt{\tau}.
\]
\end{block}
\end{frame}

\begin{frame}{The formula survives basic boundary checks}
\begin{pausebox}
Predict the limits when \(s\downarrow0\), \(s\uparrow\infty\), and
\(\tau\downarrow0\).
\end{pausebox}
\pause
\[
\begin{array}{lll}
s\downarrow0: & d_\pm\to-\infty, & C(t,s)\to0,\\
s\uparrow\infty: & d_\pm\to+\infty, & C(t,s)\sim s-Ke^{-r\tau},\\
\tau\downarrow0: & & C(t,s)\to(s-K)^+.
\end{array}
\]
Also \(0\le C(t,s)\le s\).
\begin{exampleblock}{Why check limits?}
They catch sign mistakes in discounting, \(\tau=T-t\), and \(d_\pm\).
\end{exampleblock}
\end{frame}

\begin{frame}{Put--call parity gives the put without another integral}
At maturity,
\[
(S_T-K)^+-(K-S_T)^+=S_T-K.
\]
Replication and no arbitrage imply
\[
\boxed{C(t,s)-P(t,s)=s-Ke^{-r\tau}.}
\]
\begin{pausebox}
Deduce the Black--Scholes put price.
\end{pausebox}
\pause
\[
\boxed{
P(t,s)
=Ke^{-r\tau}\Ncdf(-d_-)-s\Ncdf(-d_+).}
\]
\end{frame}

% ============================================================
\section{Greeks Are Hedge and Model Diagnostics}
% ============================================================

\begin{frame}{One density identity makes every derivative shorter}
Let
\[
C=s\Ncdf(d_+)-Ke^{-r\tau}\Ncdf(d_-),
\qquad
\phiN(x)=\frac{e^{-x^2/2}}{\sqrt{2\pi}}.
\]
\begin{pausebox}
Show that the two density terms satisfy a simple identity.
\end{pausebox}
\pause
Using \(d_-=d_+-\sigma\sqrt{\tau}\),
\[
\frac{\phiN(d_-)}{\phiN(d_+)}
=\exp\left(\frac{d_+^2-d_-^2}{2}\right)
=\frac{se^{r\tau}}{K}.
\]
Thus
\[
\boxed{s\phiN(d_+)=Ke^{-r\tau}\phiN(d_-).}
\]
\end{frame}

\begin{frame}{Delta is the stock holding in the replicating strategy}
\[
C(t,s)=s\Ncdf(d_+)-Ke^{-r\tau}\Ncdf(d_-),
\qquad
d_\pm=\frac{\log(s/K)+(r\pm\frac12\sigma^2)\tau}
{\sigma\sqrt{\tau}}.
\]
\begin{pausebox}
Starting from the displayed formula, differentiate \(C(t,s)\) with respect
to \(s\). Keep both \(d_\pm\) terms until they cancel.
\end{pausebox}
\pause
Since
\[
\partial_s d_+=\partial_s d_-=\frac{1}{s\sigma\sqrt{\tau}},
\]
\[
\begin{aligned}
\Delta=C_s
={}&\Ncdf(d_+)+s\phiN(d_+)\partial_sd_+\\
&-Ke^{-r\tau}\phiN(d_-)\partial_sd_-.
\end{aligned}
\]
The density identity gives
\[
\boxed{\Delta_{\mathrm{call}}=\Ncdf(d_+).}
\]
\end{frame}

\begin{frame}{Gamma measures how fast the hedge must move}
\begin{pausebox}
Differentiate \(\Delta=\Ncdf(d_+)\) once more with respect to \(s\).
\end{pausebox}
\pause
\[
\Gamma=C_{ss}
=\phiN(d_+)\partial_sd_+
=\boxed{\frac{\phiN(d_+)}{s\sigma\sqrt{\tau}}.}
\]
\begin{exampleblock}{Interpretation}
\(\Gamma>0\) for a call. The hedge changes most rapidly near the strike and
near maturity.
\end{exampleblock}
\begin{marketbox}
Continuous rebalancing is costless in the model. In practice, large gamma
means frequent trading, costs, and discrete-hedging error.
\end{marketbox}
\end{frame}

\begin{frame}{Vega quantifies volatility-model sensitivity}
\begin{pausebox}
Differentiate the call price with respect to \(\sigma\). Use the density
identity instead of expanding every term separately.
\end{pausebox}
\pause
The \(d_\pm\)-derivative terms simplify to
\[
\boxed{
\mathrm{Vega}=C_\sigma=s\phiN(d_+)\sqrt{\tau}.}
\]
\begin{block}{Interpretation}
Vega is positive: more volatility increases the value of a convex payoff.
\end{block}
\begin{simplificationbox}
Black--Scholes treats \(\sigma\) as known and constant. Market implied
volatility varies by strike and maturity.
\end{simplificationbox}
\end{frame}

\begin{frame}{Theta is quickest to compute from the PDE}
\begin{pausebox}
Use the PDE with the known call delta and gamma to find \(\Theta=C_t\).
\end{pausebox}
\pause
\[
\Theta
=rC-rs\Delta-\frac12\sigma^2s^2\Gamma.
\]
Substitution gives
\[
\boxed{
\Theta_{\mathrm{call}}
=-\frac{s\sigma\phiN(d_+)}{2\sqrt{\tau}}
-rKe^{-r\tau}\Ncdf(d_-).}
\]
\begin{alertblock}{Convention}
Here theta means \(\partial C/\partial t\), calendar time moving forward.
Some systems report the opposite sign or a one-day rescaling.
\end{alertblock}
\end{frame}

\begin{frame}{Rho measures sensitivity to the financing rate}
\begin{pausebox}
Differentiate \(C(t,s)\) with respect to \(r\).
\end{pausebox}
\pause
Because
\[
\partial_r d_+=\partial_r d_-=\frac{\sqrt{\tau}}{\sigma},
\]
the density contributions cancel, leaving
\[
\boxed{
\mathrm{Rho}_{\mathrm{call}}
=C_r=K\tau e^{-r\tau}\Ncdf(d_-).}
\]
\begin{simplificationbox}
The formula assumes a single deterministic rate used for discounting,
borrowing, and lending. Real funding and collateral conventions are richer.
\end{simplificationbox}
\end{frame}

\begin{frame}{Put Greeks follow from parity}
Differentiate
\[
P=C-s+Ke^{-r\tau}.
\]
\begin{pausebox}
Find the put delta, gamma, vega, theta, and rho.
\end{pausebox}
\pause
\[
\begin{aligned}
\Delta_P&=\Ncdf(d_+)-1=-\Ncdf(-d_+),\\
\Gamma_P&=\Gamma_C,\qquad \mathrm{Vega}_P=\mathrm{Vega}_C,\\
\Theta_P&=-\frac{s\sigma\phiN(d_+)}{2\sqrt{\tau}}
+rKe^{-r\tau}\Ncdf(-d_-),\\
\mathrm{Rho}_P&=-K\tau e^{-r\tau}\Ncdf(-d_-).
\end{aligned}
\]
\end{frame}

\begin{frame}{The five main Greeks answer different questions}
\small
\begin{tabular}{>{\bfseries}l l l}
\toprule
Greek & Formula for a call & Main risk question\\
\midrule
\(\Delta\) & \(\Ncdf(d_+)\) & How many shares hedge?\\
\(\Gamma\) & \(\phiN(d_+)/(s\sigma\sqrt{\tau})\) & How fast does delta change?\\
Vega & \(s\phiN(d_+)\sqrt{\tau}\) & What if volatility is wrong?\\
\(\Theta\) & \(-s\sigma\phiN(d_+)/(2\sqrt{\tau})-rKe^{-r\tau}\Ncdf(d_-)\) & What does time cost?\\
Rho & \(K\tau e^{-r\tau}\Ncdf(d_-)\) & What if rates change?\\
\bottomrule
\end{tabular}
\begin{block}{Hierarchy}
Delta is the model hedge. Gamma controls hedge instability. Vega, theta, and
rho expose parameter and carry sensitivity.
\end{block}
\end{frame}

\begin{frame}{A Greek is local, while hedging error accumulates}
With discrete trading over \(\Delta t\),
\[
\Delta C
\approx
\Delta\,\Delta S
+\Theta\,\Delta t
+\frac12\Gamma(\Delta S)^2.
\]
\begin{pausebox}
Which idealizations become most fragile for a high-gamma option?
\end{pausebox}
\pause
\begin{itemize}
  \item continuous trading;
  \item zero transaction costs and unlimited liquidity;
  \item continuous stock paths;
  \item known constant volatility.
\end{itemize}
\end{frame}

% ============================================================
\section{What the Derivation Really Established}
% ============================================================

\begin{frame}{The complete chain is one reproducible argument}
\[
\boxed{
\begin{gathered}
\text{GBM model + frictionless complete market}\\
\Downarrow\\
\text{self-financing replication with }\Delta=p_s\\
\Downarrow\\
p_t+rsp_s+\tfrac12\sigma^2s^2p_{ss}-rp=0\\
\Downarrow\\
p(t,s)=e^{-r\tau}\E[g(S_T^r)\mid S_t^r=s]\\
\Downarrow\\
C=s\Ncdf(d_+)-Ke^{-r\tau}\Ncdf(d_-).
\end{gathered}}
\]
\end{frame}

\begin{frame}{Final check: identify the decisive assumption}
\begin{pausebox}
\begin{enumerate}
  \item The \(\dd W\) exposure can be eliminated.
  \item The hedged position earns \(r\).
  \item The call has a \(d_\pm\) closed form.
  \item Delta can be followed exactly.
\end{enumerate}
\end{pausebox}
\pause
\begin{enumerate}
  \item One traded stock spans the one Brownian risk: completeness.
  \item No arbitrage and unrestricted borrowing/lending at \(r\).
  \item Constant \(r,\sigma\), lognormal dynamics, no dividends, European payoff.
  \item Continuous frictionless trading with predictable admissible strategies.
\end{enumerate}
\end{frame}

\begin{frame}{Keep these six ideas}
\begin{enumerate}
  \item It\^o's second-order term creates gamma exposure.
  \item Self-financing is a gain equation, not merely a value identity.
  \item Matching diffusion determines delta.
  \item Matching drift gives the PDE and removes \(\mu\).
  \item The logarithmic transform turns the PDE into heat flow.
  \item Verification connects the PDE solution, discounted expectation, and
        replicating strategy.
\end{enumerate}
\begin{exampleblock}{Next conceptual step}
Risk-neutral measure changes explain probabilistically why the auxiliary
drift is \(r\).
\end{exampleblock}
\end{frame}

\begin{frame}{Optional extension: continuous dividends}
Suppose the stock pays a continuous dividend yield \(q\). Under the
risk-neutral measure,
\[
\frac{\dd S_t}{S_t}=(r-q)\dd t+\sigma\dd W_t^\Q,
\qquad \dd B_t=rB_t\dd t.
\]
For a European call with payoff \((S_T-K)^+\),
\[
C(t,s)=e^{-r\tau}\E^\Q\!\left[(S_T-K)^+\mid S_t=s\right],
\qquad \tau=T-t.
\]
\begin{pausebox}
Derive the pricing PDE and the call formula. What changes in the stock term,
and what stays the same in the strike term?
\end{pausebox}
\end{frame}

\end{document}
